% Draft of the paper: the abstract is written and the sections are outlines.
%
% Build: tectonic paper/main.tex
% Style: neurips_2026.sty, from the NeurIPS 2026 formatting instructions, downloaded on 2026-10-11 from
% https://media.neurips.cc/Conferences/NeurIPS2026/Formatting_Instructions_For_NeurIPS_2026.zip
%
% The option of the style sets the version of the paper:
%   preprint              the arXiv version, with the author shown (used now)
%   sglblindworkshop      a workshop submission with the author shown
%   dblblindworkshop      an anonymous workshop submission
%   main                  an anonymous submission to the main track; add `final` for the camera-ready version
\documentclass{article}
\usepackage[preprint]{neurips_2026}

\usepackage[utf8]{inputenc}
\usepackage[T1]{fontenc}
\usepackage{hyperref}
\usepackage{url}
\usepackage{booktabs}
\usepackage{amsfonts}
\usepackage{nicefrac}
\usepackage{microtype}
\usepackage{xcolor}
\usepackage{graphicx}

% ---------------------------------------------------------------------------------------------------------------------
% Metadata
% ---------------------------------------------------------------------------------------------------------------------
% Title:           to be decided
% Author:          Rushab Munot, independent researcher, no affiliation
% Contact:         to be decided
% arXiv category:  cs.CL (to be confirmed)
% Licence:         to be decided
% Version, date:   to be decided
% Keywords:        to be decided
% Code:            see "Code and data availability" and Appendix "Code"
% ---------------------------------------------------------------------------------------------------------------------

\title{Title to be decided}

\author{%
  Rushab Munot \\
  Independent Researcher \\
  % \texttt{email to be decided} \\
}

\begin{document}

\maketitle

\begin{abstract}
Decision models return a probability distribution over a supplied list of choices. We test whether an option's reported probability can be used as the probability that the option is correct. We evaluate Jev, Liquid d1, and GPT-6 Luna Decisions on MetaTool, MMLU, HotpotQA, and ECtHR, constructing lists with one, several, or no correct choices and with five copies of one correct choice. On lists with two to four correct choices, the highest-scored correct choice receives a mean probability of 0.66 to 0.85, while the lowest receives 0.11 to 0.19. A threshold of 0.25 retains every correct choice in only 15\% to 35\% of these answers. With five semantically identical choices, Liquid d1 and GPT-6 Luna Decisions assign mean probabilities of 0.50 to 0.81 and 0.67 to 0.79 to the most favored copy, compared with an equal-share reference of 0.20. On lists with no correct choice, Jev and Liquid d1 return choices with mean top probabilities of 0.49 to 0.87 across the tested conditions. On lists with one correct choice, mean top probability differs from accuracy by at most 0.14. The results show that a choice probability describes allocation within the supplied list and cannot generally be read as an independent probability that its option is correct.
\end{abstract}

\section{Introduction}

\section{Background: decision models and the Choice question}

\section{Related work}

\section{Experimental setup}

\subsection{Models}

\subsection{Datasets}

\subsection{Lists}

\subsection{Measures}

\section{Results}

\subsection{Lists with several correct options}

\subsection{Lists of identical options}

\subsection{Lists with no correct option}

\subsection{Lists with one correct option}

\subsection{The documented example}

\section{Discussion}

\section{Limitations}

\section{Conclusion}

\section*{Code and data availability}
% Repository URL:          to be decided
% Commit of the results:   to be decided
% Licence of the code:     to be decided
% Saved model responses:   to be decided (released or not)
% Command that rebuilds the figures and tables: to be decided

\begin{ack}
% Acknowledgments, funding and the use of AI assistance. To be written.
\end{ack}

\bibliographystyle{plainnat}
\bibliography{references}

% ---------------------------------------------------------------------------------------------------------------------
\appendix

\section{Terms and measures}

\section{Models and API settings}

\section{Datasets and licences}

\section{Construction of the lists}

\section{Instructions sent to the models}

\section{Results by run}

\section{Refusals of GPT-6 Luna Decisions}

\section{Support tickets for the documented example}

\section{Examples}

\section{Cost}

\section{Code}
% Repository URL, commit, and the commands that download the data, send the requests and rebuild every figure and table.

\section{Run metadata}
% Dates of the runs, model names and versions, numbers of requests, input tokens, and the commit of each results folder.

% A submission to the main track of NeurIPS also needs the checklist of the formatting instructions:
% \input{checklist.tex}

\end{document}
